\documentclass[11pt,a4paper]{article}

% -------------------- Langue / encodage --------------------
\usepackage[T1]{fontenc}
\usepackage[french]{babel}
\usepackage{array}
\usepackage{calc}

% -------------------- Mise en page --------------------
\usepackage{geometry}
\geometry{
  a4paper,
  landscape,
  left=1.6cm,
  right=1.6cm,
  top=1.0cm,
  bottom=1.2cm
}

% -------------------- Maths / mise en forme --------------------
\usepackage{amsmath,amssymb}
\usepackage{parskip}
\usepackage[explicit]{titlesec}
\usepackage[normalem]{ulem} % pour \uline sans casser \emph
\usepackage{titlesec}

\usepackage{fancyhdr}
\usepackage{lastpage}
\usepackage[hidelinks]{hyperref}
\usepackage{enumitem}
\usepackage[table]{xcolor} % <-- IMPORTANT : pour \rowcolor dans les tableaux
\usepackage{paracol}

% -------------------- TikZ + tkz-tab --------------------
\usepackage{tikz}
\usetikzlibrary{arrows.meta,calc}
\usepackage{tkz-tab}

% -------------------- Variables --------------------
\newcommand{\FicheTitre}{Primitives et équations différentielles -- Fiche de cours}
\newcommand{\Niveau}{Terminale générale -- Mathématiques}
\newcommand{\Annee}{Année scolaire 2025/2026}

\newcommand{\SiteURL}{https://naoufel-coursparticuliers.netlify.app/}
\newcommand{\SiteTexte}{naoufel-coursparticuliers.netlify.app}

% -------------------- Footer --------------------
\pagestyle{fancy}
\fancyhf{}
\lfoot{\textit{\Niveau \;--\; \Annee}}
\cfoot{\thepage/\pageref{LastPage}}
\rfoot{\href{\SiteURL}{\texttt{\SiteTexte}}}
\renewcommand{\headrulewidth}{0pt}
\renewcommand{\footrulewidth}{0.4pt}

% -------------------- Spacing (anti débordement) --------------------
\setlength{\parskip}{4pt}
\setlength{\parindent}{0pt}
\setlist[itemize]{leftmargin=*, itemsep=2pt, topsep=2pt}

% -------------------- Titres --------------------
%\titleformat{\section}{\large\bfseries}{\thesection.}{0.5em}{} %noir
%\titleformat{\subsection}{\normalsize\bfseries}{\alph{subsection}.}{0.5em}{}\titleformat{\section}
% --- TITRES colorés + petite ligne sous le TEXTE (sans souligner le numéro) ---
% SECTION : numéro rouge (non souligné) + titre rouge souligné
\titleformat{\section}[hang]
  {\normalsize\bfseries}
  {\textcolor{red}{\thesection.}}
  {0.2em}
  {\textcolor{red}{\uline{#1}}}

% SUBSECTION : lettre verte (non soulignée) + titre vert souligné
\titleformat{\subsection}[hang]
  {\normalsize\bfseries}
  {\hspace*{1.2em}\textcolor{green!60!black}{\thesubsection}}
  {0.4em}
  {\textcolor{green!60!black}{\uline{#1}}}



% -------------------- Bandeau titre --------------------
\newcommand{\TitreBandeau}[1]{%
  \noindent\hrule
  \vspace{0.55em}
  \begin{center}
    {\LARGE\bfseries #1}
  \end{center}
  \vspace{0.55em}
  \noindent\hrule
  \vspace{0.9em}
}

% -------------------- Helpers --------------------
\newcommand{\defi}{\textbf{Définition :} }
\newcommand{\prop}{\textbf{Propriété :} }
\newcommand{\rem}{\textbf{Remarque :} }

% ============================
% Tableau 1 : Primitives usuelles (2 colonnes)
% ============================
\newcommand{\TableauPrimitivesUsuelles}{%
\begin{center}
\setlength{\tabcolsep}{10pt}
\renewcommand{\arraystretch}{1.6}
\begin{tabular}{|>{\centering\arraybackslash}m{0.43\linewidth}|>{\centering\arraybackslash}m{0.43\linewidth}|}
\hline
\rowcolor{blue!15}\textbf{Fonction $f(x)$} & \textbf{Une primitive $F(x)$} \\
\hline
$a$ (constante) & $ax$ \\
\hline
$x$ & $\dfrac{x^{2}}{2}$ \\
\hline
$x^{n}\ (n\in\mathbb{Z},\ n\neq -1)$ & $\dfrac{x^{n+1}}{n+1}$ \\
\hline
$\dfrac{1}{x}$ & $\ln(x)$ \\
\hline
$\dfrac{1}{x^{n}}\ (n\in\mathbb{Z},\ n\neq 1)$ & $-\dfrac{1}{(n-1)x^{n-1}}$ \\
\hline
$\dfrac{1}{\sqrt{x}}$ & $2\sqrt{x}$ \\
\hline
$\sin(x)$ & $-\cos(x)$ \\
\hline
$\cos(x)$ & $\sin(x)$ \\
\hline
$e^{x}$ & $e^{x}$ \\
\hline
\end{tabular}
\end{center}

}

% ============================
% Tableau 2 : Primitives de fonctions composées (2 colonnes)
% ============================
\newcommand{\TableauPrimitivesComposees}{%
\begin{center}
\setlength{\tabcolsep}{10pt}
\renewcommand{\arraystretch}{1.8}
\begin{tabular}{|>{\centering\arraybackslash}m{0.43\linewidth}|>{\centering\arraybackslash}m{0.43\linewidth}|}
\hline
\rowcolor{blue!15}\textbf{Fonction $f(x)$ de la forme} & \textbf{Une primitive $F(x)$} \\
\hline

$u'\,u^{n}\ (n\neq -1;0)$ & $\dfrac{u^{n+1}}{n+1}$ \\
\hline
$\dfrac{u'}{u}$ & $\ln|u|$ \\
\hline
$\dfrac{u'}{u^{2}}$ & $-\dfrac{1}{u}$ \\
\hline
$\dfrac{u'}{\sqrt{u}}$ & $2\sqrt{u}$ \\
\hline
$u'\,e^{u}$ & $e^{u}$ \\
\hline
$u'\,\cos(u)$ & $\sin(u)$ \\
\hline
$u'\,\sin(u)$ & $-\cos(u)$ \\
\hline
$(v'\!\circ u)\,u'$ & $v\!\circ u$ \\


\hline
\end{tabular}
\end{center}

}

\begin{document}

\TitreBandeau{\FicheTitre}

% -------------------- 2 colonnes + ligne verticale --------------------
\columnratio{0.5,0.5}
\setlength{\columnsep}{1.2cm}
\setlength{\columnseprule}{0.6pt}
\def\columnseprulecolor{\color{black}}

\begin{paracol}{2}

% ====== COLONNE GAUCHE ======
\section{Équation différentielle $y'=f$}

\defi Une équation différentielle est une équation dans laquelle
l'inconnue est une fonction et où intervient au moins une de ses dérivées.

\subsection{Primitive d'une fonction continue sur un intervalle}
\defi On dit que $F$ est une \textbf{primitive} de $f$ sur un intervalle $I$
lorsque $F$ est dérivable sur $I$ et
\[
F'(x)=f(x)\quad \text{pour tout } x\in I.
\]
\prop Toute fonction \textbf{continue} sur un intervalle admet des primitives sur cet intervalle
(\textit{à une constante près}).

\subsection{Primitives usuelles (référence)}
\TableauPrimitivesUsuelles

% >>> Basculer sur la colonne de droite <<<
\switchcolumn
\setcounter{section}{1}
\setcounter{subsection}{2}

\subsection{Primitives de fonctions composées}
\TableauPrimitivesComposees



% ====== COLONNE DROITE ======
\setcounter{section}{1}

\section{Équation différentielle $y'=ay$}

\prop L'équation différentielle $y'=ay$ est une équation différentielle linéaire homogène
du premier ordre (sans second membre).

Les solutions sont de la forme :
\[
y(x)=Ce^{ax}\qquad (C\in\mathbb{R})
\]
\rem Pour déterminer $C$, on utilise une condition (valeur initiale, point imposé, etc.).

% >>> Basculer sur la colonne de droite <<<
\switchcolumn
\setcounter{section}{2}


\section{Équation différentielle $y'=ay+b\ (a\neq 0)$}
Les solutions de l'équation différentielle $y'=ay+b$ sont de la forme :
\[
y(x)=Ce^{ax}-\frac{b}{a}\qquad (C\in\mathbb{R})
\]
\rem Pour déterminer $C$, on utilise une condition.

\section{Équation différentielle $y'=ay+f\ \ (a\neq 0)$}
Les solutions de l'équation différentielle $y'=ay+f(x)$ sont de la forme :
\[
y(x)=Ce^{ax}+\varphi(x)\qquad (C\in\mathbb{R})
\]
où $\varphi(x)$ est une fonction de même nature que $f(x)$ (solution particulière adaptée à $f$).

\rem Pour déterminer $C$, on utilise une condition.

\section{Principe de résolution des équations différentielles}
\prop
\[
\text{Solution générale}=\text{Solution homogène}+\text{Solution particulière}.
\]
\rem On commence par résoudre l'équation homogène associée, puis on ajoute une solution particulière
du second membre (quand il y en a un).

\end{paracol}
\end{document}
